\section{Introduction}
\label{sec:introduction}

The rapid expansion of photovoltaic generation has increased the importance of
solar-resource forecasts for system planning, energy management, and grid
integration. Forecast usefulness is nevertheless conditional on temporal
resolution, prediction horizon, available covariates, and evaluation protocol
\cite{antonanzas2016,voyant2017,yang2018}. A model that performs well for an
hourly three-day trajectory does not necessarily remain competitive for daily
or monthly planning.

Industrial decarbonization makes this distinction especially relevant for
manufacturing facilities. Electric-vehicle factories provide a geographically
diverse case study because their locations span different latitudes, seasonal
cycles, and atmospheric regimes. In this work, the factories identify the
geographic points at which solar irradiance is studied; no production,
consumption, or proprietary operational data are used.

Two preceding studies motivate this journal extension. The first configured a
compact TimesNet for direct 72-hour global horizontal irradiance (GHI)
forecasting and evaluated the 24-, 48-, and 72-hour prefixes. The second
evaluated a Dilated Recurrent Neural Network (DilatedRNN) for one-month-ahead
forecasting from short monthly series. Their results cannot yet establish which
architecture is preferable across scales because the targets, context lengths,
forecast origins, competitors, and primary metrics differ.

This extension therefore has three objectives: (i) provide a unified and
reproducible description of the data, transformations, architectures, and
evaluation rules; (ii) evaluate TimesNet and DilatedRNN at shared hourly,
daily, and longer planning horizons; and (iii) characterize when performance
changes with location, horizon, and climate-related variability. Its intended
contribution is a matched multi-resolution comparison, rather than a direct
juxtaposition of incomparable errors from the preceding manuscripts.

The remainder of the article reviews the theoretical background, documents the
data and feature engineering, defines the common experimental protocol,
presents the evaluation metrics separately, and analyzes aggregate and adverse
cases. Full hyperparameter grids and location-level results are assigned to the
appendix.

